\heading{What the deterministic offset actually measures}
\label{sec:physics-offset}
The colored curve in v5.8.2 Figure~3 is $a(m)=r(B;m)$: the signed signal-fit root obtained when the input histogram is replaced by the smooth expected counts $B$, with no Poisson draw. It is not the observed excess, an estimated number of signal events, or a fraction of signal removed. There are two distinct fits. First, the \emph{generating source} $B$ is a GP fit through every observed bin, with its kernel anchored at 76~MeV. Second, at each tested mass, the ordinary extraction excludes $\pm2.25\sigma_m$, predicts inside that window from the sidebands, and profiles the signal and background nuisance parameters. The source fit has no moving exclusion window. The extraction has mass-dependent reviewed kernels and count-dependent conditioning. These two operations need not reproduce one another exactly: a smooth $B$ can produce a nonzero signed signal estimate when it is rescanned.

A positive $a$ means that this particular noiseless reference looks signal-like to the extraction at that mass; a negative $a$ means a deficit relative to the extraction's interpolated background. Curvature, interpolation, the different source/extraction kernels and nuisance treatment can contribute. This diagnoses the response to the chosen $B$; it does not identify which component of the measured spectrum is continuum or signal. The unchanged observed root remains available separately.

The other two curves in the old Figure~3 are already transformed quantities. Its black curve is $\operatorname{mean}[(r_{\rm toy}-a)/s]$; its gray dashed curve is $\operatorname{SD}[(r_{\rm toy}-a)/s]$, where $s$ is the linear-response fluctuation width. They should lie near zero and one if the conditional approximation works. They are neither raw toy roots nor additional offsets to subtract. In particular, a black curve near zero does not imply $a=0$. The nonlinearity of the scan also means that $r(B)$ need not equal the exact toy mean, although their agreement can be checked.

\fig{physics_fig3_explained}{A clearer rendering of the 2016 part of v5.8.2 Figure~3. Panel A adds the unchanged observed scan; panel B compares the raw toy mean with the deterministic offset. Panels C and D then show the centered mean and standardized width separately. Gray bands mark only approximate pointwise one-standard-error scales for 256 independent Gaussian draws, $1/\sqrt{256}$ and $1/\sqrt{2(255)}$; they are not simultaneous acceptance bands. The mass coordinates are correlated, and all curves condition on the same fixed source.}

For example, at 90.5~MeV the stored values are $a=0.71285$ and $s=0.98203$. The 256 raw toy roots have mean $0.65653$ and SD $1.04126$. Their centered mean is about $-0.057$, explaining how the old figure can show a sizeable colored offset together with a black curve near zero. At 76~MeV the nominal-GP offset is $-0.40839$; the much larger archived-stress response belongs to a different reference. Neither value alone proves that a particle signal was fitted away.

\heading{Source learning and signal recovery are different tests}
The signal-learning concern is nevertheless supported by a separate experiment. In v5.8.1, a signal was added to the \emph{source data}, the unmasked GP source was rebuilt, and the change in $B$ was measured. The nominal 2016 source absorbed 88.7\% and 88.4\% of the added counts inside the 76 and 90~MeV windows, respectively, and about 72\% of the signal-shaped component. By contrast, adding signal to an already frozen nominal source and running the masked extraction recovered roughly 96\% of the injected amplitude. Successful extraction recovery does not qualify a source that is subsequently rebuilt through the signal. The absorbed fraction is a deterministic transfer coefficient, not a detection efficiency, false-positive rate or a statement that this fraction of the actual observed excess is signal.

The practical concern is therefore the full procedure $\text{data}\to\widehat B\to\text{reference calibration}$, not merely the masked extraction. A learned feature can enter the putative null, affect its deterministic response and covariance, and change the resulting score. The source estimate was held fixed in the v5.8.2 toys, so their agreement tests propagation \emph{conditional on that estimate}; it does not test repeated source learning or establish a background-only physical truth. A smaller stress RMS is useful for diagnosing compatibility between two modeling operations, but cannot choose the background model by itself.

We retain $\pm2.25\sigma_m$ throughout. Widening only the extraction mask to $3\sigma_m$ would leave the unmasked source fit untouched, so it would not cure this source-learning mechanism. For an ideal Gaussian template, the full tails beyond these half-widths are about 2.445\% and 0.270\%, respectively; those template-tail fractions are distinct from the much larger source-absorption numbers. No new width scan was performed.

\heading{The proposed high-psum 1\% source: a bounded test}
\label{sec:physics-onepercent}
A real candidate is available: the 2021 file \path{final_1pct_invM_v7.root}, histogram \path{preselection/h_invM_psumgt2p8_8000}. Its key explicitly labels $p_{\rm sum}>2.8$; the stored histogram contains 6,336,937 visible counts. The existing selection audit finds no event identities, run list, prompt-TC certificate or exact luminosity ratio. Its counts are 4.476\% of the native-10\% histogram overall, and 5.148\% in 87--97~MeV. These are count ratios, not measured exposure ratios. The native-10\% input comes from the documented prompt-TC production; the high-psum key alone does not establish matching reconstruction or all other cuts. This is a 2021 control and cannot supply the full-2016 offset without a separate, validated cross-dataset transfer model.

We tested the proposal as an explicitly conditional calculation. The high-psum histogram was exactly rebinned to the current 422-bin support, smoothed using the existing 2021 source kernel at 76~MeV, and multiplied by ten. The same current extraction, resolution and $\pm2.25\sigma_m$ masks were then used at 201 integer masses from 50 to 250~MeV. The three sources are the native high-psum GP mean, that identical mean multiplied by ten, and the saved native-10\% GP mean. This is a common-method comparison, not a validation of the nominal exposure transfer. No kernel was selected to improve these results.

\fig[.90]{physics_source_transfer}{New fixed-width 2021 controls. A: the nominally scaled high-psum source is only about 13--52\% of the native-10\% GP mean over 50--250~MeV, with substantial shape dependence. B: deterministic source responses using the same extraction settings. C: fixed-source recovery versus source rebuilding after a known signal injection. The two source-learning bars use different projections of the same rebuilt mean; neither is a probability. The smaller offsets in B do not repair the transfer mismatch or the absorption in C.}

The deterministic RMS offsets are $0.0554$, $0.0283$ and $0.2576$ for those three sources, respectively. Offset scaling is not necessarily $\sqrt{10}$: reconditioning the extraction at the new count level changes its GP noise and interpolation. We also generated 64 complete independent Poisson spectra for each target-scale source and fitted six fixed anchors, 65, 76, 90, 92, 120 and 180~MeV, for 768 fits. The same draw is reused at all six masses. Centered means span $-0.188$ to $0.187$, and standardized widths span $0.914$ to $1.080$, consistent with only a modest moment check. These six-anchor toys are not a full-search calibration or a rare-tail measurement; their approximate mean and width Monte Carlo errors are 0.125 and 0.089. All 603 deterministic and 768 toy fits satisfy the retained $2\times10^{-7}$ score criterion.

An additional injection test addresses the key physical issue. At 76 and 90~MeV, add a five-error target-scale Gaussian signal to the fixed $10B_{1\%}$ spectrum. The extraction recovers 96.67\% and 96.68\% of its amplitude. Separately, add one tenth of that same signal to the \emph{1\% source counts}, rebuild the unmasked GP, and multiply the new source by ten. The source absorbs 84.04\% and 77.83\% of the target window yield, or 69.02\% and 63.84\% of its Poisson-weighted matched shape. Thus a smaller, statistically separate sample would help limit reuse of the search fluctuations, but a genuine signal present in both samples can still be learned as background.

\heading{What would make a 1\%-to-10\% calibration defensible?}
The useful next experiment is a prediction from an audited control sample, with its uncertainty, rather than a fixed smoothed histogram treated as exact. First verify the event selection, processing version, exposure and event overlap. If the 1\% is included in the 10\%, either reserve a disjoint search sample or model their joint sampling; the two are not independent merely because the toy random seeds differ. Freeze the source construction, search statistic and mass grid before inspecting the validation sample. Assess continuum interpolation and source-injection absorption across the declared domain, including boundaries, and retain a signal-aware source model or controls where absorption is acceptably bounded.

Source uncertainty matters especially when increasing the exposure. If a disjoint source has mean counts $\lambda$ and the target has mean $k\lambda$, the deliberately simple unsmoothed example $N_{\rm target}-kN_{\rm source}$ has variance $(k+k^2)\lambda$, whereas treating the source as fixed keeps only $k\lambda$. A nested source instead contributes a covariance term. This example is not the variance of the actual GP statistic: smoothing and the extraction response must propagate the relevant modes, or an outer ensemble must repeat source estimation and target generation. For each specified underlying background truth, jointly generate control and target samples with the correct exposures and event overlap. Refit the source from that control realization, apply the declared transfer and calibration, and scan the corresponding target realization. Repeat every data-dependent choice that would be made in the real analysis. Generating the target from each noisy fitted source instead would conceal the source-estimation error that this test is intended to expose. Use separate realizations to validate local null tails, signal recovery and the maximum distribution. The existing fixed-source GP/Poisson agreement is a valuable inner check within this larger procedure.

The present high-psum calculation therefore supports keeping this source as a sensitivity control and identifies what must be qualified next. It does not justify replacing the released local significances with a new offset, certifying source independence, or adopting the source merely because its deterministic response is small.
